% facitliste.tex — Facitliste
% Hamilton-Jacobi-kompendiet

\chapter*{Facitliste}
\addcontentsline{toc}{chapter}{Facitliste}
\markboth{FACITLISTE}{FACITLISTE}

Denne facitliste giver facit og korte løsningsskitser til
kompendiets opgaver — ikke fulde udførte udregninger (de fleste
opgaver er netop bedt om at \emph{gennemføre} en udregning, hvis
opskrift allerede står i teksten). Brug den til at tjekke dit eget
resultat og din metode, ikke som erstatning for selv at regne.

\section*{Kapitel 1: Newtonsk mekanik og variationsprincipper}
\addcontentsline{toc}{section}{Kapitel 1}

\paragraph{Opgave 1.} Newtons anden lov langs banen og vinkelret på
den giver to ligninger: $m\ddot s = mg\sin\alpha$ (langs skråplanet)
og $N=mg\cos\alpha$ (vinkelret, normalkraften). $N$ indgår slet ikke
i den første ligning — den ``falder ud'' ved, at normalkraften per
definition ikke har nogen komponent langs den tilladte
bevægelsesretning, ganske som i D'Alembert-formuleringen, hvor
tvangskraften forsvinder, fordi $\vb{N}\cdot\delta\vb{r}=0$ for enhver
tilladt virtuel forskydning. Facit: $\ddot s=g\sin\alpha$.

\paragraph{Opgave 2.} Med $\theta$ som generaliseret koordinat er
$\vb{r}=R(\sin\theta,-\cos\theta)$ (målt fra ringens centrum, med
$\theta=0$ i bunden), så $\delta\vb{r}=R(\cos\theta,\sin\theta)\,\delta\theta$.
Tyngdekraften er $\vb{F}=(0,-mg)$, og tvangskraften (ringens
normalkraft) står vinkelret på $\delta\vb{r}$ og bidrager derfor ikke.
D'Alemberts princip, $(\vb{F}-m\ddot{\vb{r}})\cdot\delta\vb{r}=0$, giver
efter udregning af $\ddot{\vb{r}}$ facit
\begin{equation*}
  \ddot\theta = -\frac{g}{R}\sin\theta,
\end{equation*}
den velkendte pendulligning (her for en perle på en ring i stedet for
en snor, men matematisk identisk).

\section*{Kapitel 2: Lagrange-formalisme}
\addcontentsline{toc}{section}{Kapitel 2}

\paragraph{Opgave 1.} Direkte udregning af $x,y,z$ for et punkt på
den roterende ring (radius $R$, vinkelhastighed $\omega$ om den
lodrette akse, $\theta$ målt fra symmetriaksen) og differentiation
giver, efter sammentrækning af $\sin^2+\cos^2=1$-led, netop
$T=\frac12mR^2\dot\theta^2+\frac12mR^2\omega^2\sin^2\theta$ — første
led er den ``sædvanlige'' kinetiske energi fra bevægelse langs
ringen, andet led kommer fra rotationen om aksen.

\paragraph{Opgave 2.} $L=\frac12m(\dot r^2+r^2\dot\phi^2)-V(r)$.
$\partial L/\partial\phi=0$, så $\phi$ er cyklisk og
$p_\phi=mr^2\dot\phi$ er bevaret. Euler-Lagrange-ligningen for $r$,
$m\ddot r = mr\dot\phi^2 - V'(r)$, bliver efter substitution af
$\dot\phi=p_\phi/(mr^2)$:
\begin{equation*}
  m\ddot r = \frac{p_\phi^2}{mr^3} - V'(r) = -V_\text{eff}'(r),
  \qquad V_\text{eff}(r)=V(r)+\frac{p_\phi^2}{2mr^2}.
\end{equation*}
Bevægelsen i $r$ er altså en endimensional bevægelse i det effektive
potential $V_\text{eff}$, hvor centrifugalleddet $p_\phi^2/(2mr^2)$
optræder som et ekstra frastødende bidrag.

\paragraph{Opgave 3.} Er $\vb{r}_\alpha=\vb{r}_\alpha(q)$ uden
eksplicit tidsafhængighed, er $\dot{\vb{r}}_\alpha=\sum_j(\partial\vb{r}_\alpha/\partial q_j)\dot q_j$
lineær i $\dot q_j$, så $T=\frac12\sum_\alpha m_\alpha\dot{\vb{r}}_\alpha^2$
er en \emph{ren} andengradsform i $\dot q_j$ (ingen lineære eller
konstante led) — en homogen funktion af grad $2$. Eulers sætning giver
$\sum_i\dot q_i\,\partial T/\partial\dot q_i=2T$. Da $V$ ikke afhænger
af $\dot q_i$, er $\partial L/\partial\dot q_i=\partial T/\partial\dot q_i$,
så $H=\sum_i\dot q_ip_i-L=2T-(T-V)=T+V$.

\section*{Kapitel 3: Hamilton-formalisme}
\addcontentsline{toc}{section}{Kapitel 3}

\paragraph{Opgave 1.} Med $L=\frac12m(\dot r^2+r^2\dot\phi^2)-V(r)$ er
$p_r=m\dot r$, $p_\phi=mr^2\dot\phi$. Legendre-transformationen giver
\begin{equation*}
  H = p_r\dot r+p_\phi\dot\phi-L = \frac{p_r^2}{2m}+\frac{p_\phi^2}{2mr^2}+V(r),
\end{equation*}
og Hamiltons ligninger $\dot r=p_r/m$, $\dot p_r=p_\phi^2/(mr^3)-V'(r)$,
$\dot\phi=p_\phi/(mr^2)$, $\dot p_\phi=0$ er ækvivalente med
Euler-Lagrange-resultaterne fra kapitel 2, opgave 2 (elimineres
$p_\phi=mr^2\dot\phi$, genfindes $V_\text{eff}$-ligningen).

\paragraph{Opgave 2.} $p_i=\partial L/\partial\dot r_i=m\dot r_i+eA_i$,
altså $\vb{p}=m\dot{\vb{r}}+e\vb{A}$. Legendre-transformationen giver,
efter at have løst $\dot{\vb{r}}=(\vb{p}-e\vb{A})/m$ og indsat,
\begin{equation*}
  H = \frac{(\vb{p}-e\vb{A})^2}{2m} + e\Phi.
\end{equation*}
Her er $H=T+V=\frac12m\dot{\vb{r}}^2+e\Phi$ (den mekaniske energi),
men udtrykt ved $\vb{p}$ i stedet for $\dot{\vb{r}}$ —$H$ selv er
altså stadig den mekaniske energi, men den kanoniske impuls $\vb{p}$
er \emph{ikke} $m\dot{\vb{r}}$, hvilket er pointen: $H$ og $T+V$ er
samme \emph{tal}, men som \emph{funktioner} af de kanoniske variable
er $H(\vb{r},\vb{p})\neq T(\vb{r},\dot{\vb{r}})+V(\vb{r})$, når man
udtrykker dem korrekt.

\paragraph{Opgave 3.} For $E$ under potentialbrøndens top er banerne
lukkede ellipselignende kurver om oprindelsen (libration) — ganske
som for den harmoniske oscillator, men fladet ud/deformeret nær
$|x|=a$, hvor den uendelige væg tvinger $p$ til at ``springe''
(fortegnsskift) i stedet for at aftage glat. Når $E$ nærmer sig
energien for $|x|=a$ (dvs.\ oscillatorens egen energi ved amplitude
$a$), nærmer faserumsbanen sig en almindelig oscillator-ellipse
\emph{indeni} boksen, uberørt af væggene; \emph{over} den energi er
der ingen tilladte baner (partiklen kan ikke være der).

\section*{Kapitel 4: Poisson-parenteser og kanonisk struktur}
\addcontentsline{toc}{section}{Kapitel 4}

\paragraph{Opgave 1.} Fra definitionen
$\{f,g\}=\sum_i(\partial f/\partial q_i\,\partial g/\partial p_i-\partial f/\partial p_i\,\partial g/\partial q_i)$
følger antisymmetri, $\{f,g\}=-\{g,f\}$, direkte ved at bytte $f$ og
$g$ i hvert led. Leibniz-reglen følger af produktreglen anvendt på
$\partial(fg)/\partial q_i=f\,\partial g/\partial q_i+g\,\partial f/\partial q_i$
(og tilsvarende for $\partial/\partial p_i$), indsat i definitionen og
omgrupperet i led med $f$ udenfor parentesen og led med $g$ udenfor.

\paragraph{Opgave 2.} Jacobi-identiteten,
$\{\{f,g\},H\}+\{\{g,H\},f\}+\{\{H,f\},g\}=0$, omskrives med
antisymmetri til $\{\{f,g\},H\}=-\{\{g,H\},f\}-\{\{H,f\},g\}=\{f,\{g,H\}\}-\{g,\{f,H\}\}$.
Er $\{f,H\}=0$ og $\{g,H\}=0$, forsvinder begge led på højresiden, så
$\{\{f,g\},H\}=0$ — $\{f,g\}$ er selv bevaret (forudsat
$\partial f/\partial t=\partial g/\partial t=0$, så det totale
tidsafledte reducerer til Poisson-parentesen med $H$ alene).

\paragraph{Opgave 3.} Direkte udregning (med
$L_x=yp_z-zp_y$, $L_y=zp_x-xp_z$) led for led i definitionen giver
$\{L_x,L_y\}=L_z$ (og cyklisk $\{L_y,L_z\}=L_x$, $\{L_z,L_x\}=L_y$).
Tilsvarende: $\{L_z,x\}=\{xp_y-yp_x,x\}=-y\{p_x,x\}\cdot(-1)=-y$ (kun
leddet med $p_x$ og $x$ bidrager, resten er nul), og
$\{L_z,y\}=x$ ved samme regning. Dette er netop formen af en
infinitesimal rotation om $z$-aksen, $\delta x=-\varepsilon y$,
$\delta y=\varepsilon x$, jf.\ $\delta f=\varepsilon\{f,L_z\}$.

\section*{Kapitel 5: Kanoniske transformationer}
\addcontentsline{toc}{section}{Kapitel 5}

\paragraph{Opgave 1.} Variation af $S$ med $\delta q_i(t_1)=\delta q_i(t_2)=0$
men $\delta p_i$ fri giver, efter partiel integration af
$p_i\delta\dot q_i$-leddet i tid (randled forsvinder pga.\
$\delta q_i=0$ ved endepunkterne):
\begin{equation*}
  \delta S = \int dt\sum_i\left[\left(-\dot p_i-\pdv{H}{q_i}\right)\delta q_i
    + \left(\dot q_i-\pdv{H}{p_i}\right)\delta p_i\right].
\end{equation*}
Da $\delta q_i,\delta p_i$ er uafhængige og arbitrære, må hver
koefficient forsvinde for sig — netop Hamiltons ligninger.

\paragraph{Opgave 2.} Totalt differentieret, $dF_2=dF_1+\sum_i(Q_idP_i+P_idQ_i)$.
Indsættes $dF_1=\sum_i(p_idq_i-P_idQ_i)+(K-H)dt$ (fra
\eqref{eq:genererende-betingelse}), ophæves $-P_idQ_i$ og $+P_idQ_i$,
og man står tilbage med $dF_2=\sum_i(p_idq_i+Q_idP_i)+(K-H)dt$ —
netop $F_2$-relationerne $p_i=\partial F_2/\partial q_i$,
$Q_i=\partial F_2/\partial P_i$.

\paragraph{Opgave 3.} Med $Q_i=p_i$, $P_i=-q_i$:
$\{Q_i,P_j\}=\{p_i,-q_j\}=-\{p_i,q_j\}=\delta_{ij}$ (brug
$\{p_i,q_j\}=-\delta_{ij}$), så betingelsen er opfyldt. Uden
minustegnet, $Q_i=p_i,P_i=q_i$: $\{Q_i,P_j\}=\{p_i,q_j\}=-\delta_{ij}\neq\delta_{ij}$
— transformationen er \emph{ikke} kanonisk uden fortegnsskiftet
(kun ombytning af $q$ og $p$ uden fortegn bryder den symplektiske
struktur).

\paragraph{Opgave 4.} Med $G=H$, $\varepsilon=dt$:
$\delta q_i=dt\,\{q_i,H\}=dt\,\partial H/\partial p_i=\dot q_i\,dt$,
og $\delta p_i=dt\,\{p_i,H\}=-dt\,\partial H/\partial q_i=\dot p_i\,dt$
— netop Hamiltons ligninger, nu genlæst som udsagnet om, at den
infinitesimale ændring i $(q_i,p_i)$ over tidsintervallet $dt$ er
præcis den kanoniske transformation, $H$ selv genererer. Dette
bekræfter, at tidsudviklingen fra $t$ til $t+dt$ er en kanonisk
transformation (og ved sammensætning, at den endelige tidsudvikling
fra $t_1$ til $t_2$ også er det).

\section*{Kapitel 6: Hamilton-Jacobi-teori}
\addcontentsline{toc}{section}{Kapitel 6}

\paragraph{Opgave 1.} HJ-ligningen er
$\frac{1}{2m}(\partial W/\partial x)^2=E$, så
$\partial W/\partial x=\sqrt{2mE}=p_0$ (konstant), $W=p_0x$. Løsningen
$\beta=\partial W/\partial E-t=mx/p_0-t$ (konstant) giver, løst for
$x$: $x(t)=x_0+\frac{p_0}{m}t$, som forventet for en fri partikel.

\paragraph{Opgave 2.} HJ-ligningen er
$\frac{1}{2m}(\partial W/\partial x)^2+mgx=E$, så
$\partial W/\partial x=\sqrt{2m(E-mgx)}$, og
$W=-\frac{1}{3mg}\bigl(2m(E-mgx)\bigr)^{3/2}\big/m$ (ved direkte
integration). Proceduren
$\beta=\partial W/\partial E-t=\text{konst.}$ giver, efter
udregning og genindsættelse af $E=\frac12mv_0^2$ (ved $x=x_0=0$, $t=0$):
$x(t)=v_0t-\frac12gt^2$ (med $x_0=0$; generelt $x_0+v_0t-\frac12gt^2$).

\paragraph{Opgave 3.} Direkte differentiation af
$x(t)=a\sin(\omega(t+\beta))$ giver
$\dot x=a\omega\cos(\omega(t+\beta))$, så
$p=m\dot x=ma\omega\cos(\omega(t+\beta))$. Sammenholdt med
$p=\sqrt{2mE-mkx^2}=\sqrt{2mE-mka^2\sin^2(\omega(t+\beta))}$, og brug
af $E=\frac12ka^2$ (fra amplituden) og $\omega=\sqrt{k/m}$, reducerer
begge udtryk til samme $ma\omega\cos(\omega(t+\beta))$ ved brug af
$\cos^2+\sin^2=1$ — identiteten $H=E$ er opfyldt for alle $t$, netop
fordi $E$ er en bevægelseskonstant.

\section*{Kapitel 7: Virknings-vinkel-variable}
\addcontentsline{toc}{section}{Kapitel 7}

\paragraph{Opgave 1.} $p_\phi$ er konstant ($=\sqrt{2mR^2E}$, fra
$H=p_\phi^2/(2mR^2)=E$), så $J=\oint p_\phi\,d\phi=2\pi p_\phi=2\pi\sqrt{2mR^2E}$.
Løst for $E$: $E(J)=J^2/(8\pi^2mR^2)$, så
$\nu=dE/dJ=J/(4\pi^2mR^2)=p_\phi/(2\pi mR^2)$. Med $p_\phi=mRv$
($v$=banehastighed) giver dette $\nu=v/(2\pi R)$ — netop
omløbsfrekvensen.

\paragraph{Opgave 2.} Kædereglen giver
$\dot w=\partial w/\partial q\cdot\dot q=(\partial^2W/\partial q\partial J)\dot q$.
Integreres dette over én periode $\tau$ i $q$ (dvs.\ $\dot q\,dt=dq$),
fås $\Delta w=\oint\partial^2W/\partial q\partial J\,dq
=\partial/\partial J\oint\partial W/\partial q\,dq=\partial J/\partial J=1$,
sammenholdt med $\Delta w=\nu\tau$ fra \eqref{eq:hamilton-vinkel}
integreret over $\tau$, giver dette $\nu\tau=1$ — samme resultat som
\eqref{eq:delta-w-1}, nu udledt ved direkte integration af $\dot w$ i
stedet for at bytte om på differentiation og integration.

\paragraph{Opgave 3.} Med $u=1/r$ bliver
$2mE+2mk/r-\ell^2/r^2=-\ell^2u^2+2mku+2mE$, et andengradspolynomium i
$u$ med rødder $u_{1,2}=\bigl(mk\mp\sqrt{m^2k^2+2mE\ell^2}\bigr)/\ell^2$
(for $E<0$: $u_1<0<u_2$, svarende til $r_{\max}=1/u_1$ negativ i
$u$, dvs.\ reelt $r_{\max}=-1/u_1$ — vær opmærksom på fortegn ved
substitutionen $dr=-du/u^2$, som vender integrationsretningen).
Standardresultatet for $\oint\sqrt{(u_2-u)(u-u_1)}\,du$ (over ét
fuldt omløb af $u$ mellem rødderne) fører, efter tilbagesubstitution
og forenkling (brug $u_2-u_1=2\sqrt{m^2k^2+2mE\ell^2}/\ell^2$ og
$u_1u_2=-2mE/\ell^2$), frem til
$J_r=-2\pi\ell+2\pi k\sqrt{m/(-2E)}$ — facit \eqref{eq:kepler-Jr}.
(Numerisk verifikation for selvvalgte $m,k,E,\ell$ anbefales som
kontrol af mellemregningen.)

\section*{Kapitel 8: Adiabatiske invarianter}
\addcontentsline{toc}{section}{Kapitel 8}

\paragraph{Opgave 1.} $E/\omega=$ konstant og
$E=\frac12m\ell^2\omega^2\theta_0^2$ giver
$m\ell^2\omega\theta_0^2=\text{konst.}$ Med $\omega=\sqrt{g/\ell}$
($\omega\propto\ell^{-1/2}$) fås $\ell^{2}\ell^{-1/2}\theta_0^2=\ell^{3/2}\theta_0^2=\text{konst.}$,
altså
\begin{equation*}
  \theta_0 \propto \ell^{-3/4}.
\end{equation*}
Et pendul, hvis snor langsomt forkortes, svinger altså med \emph{voksende}
udslagsvinkel $\theta_0$ (selvom den fysiske amplitude $\ell\theta_0\propto\ell^{1/4}$
faktisk aftager svagt).

\paragraph{Opgave 2.} $k_n=n\pi/L$ og $E_n=\hbar^2k_n^2/(2m)=\hbar^2n^2\pi^2/(2mL^2)$.
Med $h=2\pi\hbar$, dvs.\ $\hbar=h/(2\pi)$:
$E_n=\bigl(h/(2\pi)\bigr)^2n^2\pi^2/(2mL^2)=h^2n^2/(8mL^2)$ — identisk
med Bohr-Sommerfeld-resultatet $E_n=n^2h^2/(8mL^2)$.

\section*{Kapitel 9: Klassisk feltteori}
\addcontentsline{toc}{section}{Kapitel 9}

\paragraph{Opgave 1.} For den transversale kæde er udsvinget
$\varphi_i$ vinkelret på kædens retning; hver fjeder strækkes med
$\Delta\ell\approx(\varphi_{i+1}-\varphi_i)^2/(2a)$ til laveste orden
(geometrisk, for lille udsving), så den potentielle energi i
fjederen med spænding (kraft) $\tau_0$ er
$\tau_0\Delta\ell\approx\frac{\tau_0}{2a}(\varphi_{i+1}-\varphi_i)^2$.
Dette giver, i kontinuumsgrænsen ($a\to0$, $\tau_0$ fast — modsat den
longitudinale kæde, hvor det var $ka$, der holdtes fast), samme
$\mathcal{L}=\frac12\mu\dot\varphi^2-\frac12\tau_0(\varphi')^2$, med
$\tau=\tau_0$ direkte som den fysiske snorspænding.

\paragraph{Opgave 2.} Med
$\mathcal{L}_\text{KG}=\frac12(\partial_\mu\varphi\partial^\mu\varphi-m^2\varphi^2)$
er $\partial\mathcal{L}/\partial(\partial_\mu\varphi)=\partial^\mu\varphi$
og $\partial\mathcal{L}/\partial\varphi=-m^2\varphi$. Den relativistiske
Euler-Lagrange-ligning $\partial_\mu(\partial\mathcal{L}/\partial(\partial_\mu\varphi))-\partial\mathcal{L}/\partial\varphi=0$
giver $\partial_\mu\partial^\mu\varphi+m^2\varphi=0$, dvs.\
$(\Box+m^2)\varphi=0$ med $\Box\equiv\partial_\mu\partial^\mu=\partial_t^2/c^2-\nabla^2$
— Klein-Gordon-ligningen.

\paragraph{Opgave 3.} Euler-Lagrange-ligningen for $\psi$ giver
$\partial\mathcal{L}/\partial\psi=\dot{}(\partial\mathcal{L}/\partial\dot\psi)+(\,\cdot\,)'(\partial\mathcal{L}/\partial\psi')$
(og tilsvarende for $\psi^\ast$), hvilket indsat i
$\partial\mathcal{L}/\partial\psi\cdot\delta\psi$ omskriver leddet til
en ren tidsafledt plus et led, der kan kombineres med
$\partial\mathcal{L}/\partial\dot\psi\cdot\delta\dot\psi$ til
$d/dt(\partial\mathcal{L}/\partial\dot\psi\cdot\delta\psi)$ — analogt
med udledningen af \eqref{eq:kontinuitetsligning}. Samme manøvre for
det rumlige led (med $\partial\mathcal{L}/\partial\psi'$) giver den
rumlige total-afledte. Sammenlagt for begge felter ($\psi$ og
$\psi^\ast$) fås netop \eqref{eq:ladnings-kontinuitet}.

\paragraph{Opgave 4.} Med $\psi=\varphi_1+i\varphi_2$:
$\dot\psi^\ast\dot\psi=\dot\varphi_1^2+\dot\varphi_2^2$ og
$\psi^{\ast\prime}\psi'=\varphi_1'^2+\varphi_2'^2$ (krydsledene med
$i$ ophæver hinanden i produktet af komplekst konjugerede), så
$\mathcal{L}=\mu(\dot\varphi_1^2+\dot\varphi_2^2)-\tau(\varphi_1'^2+\varphi_2'^2)$
— to uafhængige reelle strenge, hver med faktor $2$ i forhold til
\eqref{eq:L-tætthed-streng} (svarende til to frihedsgrader). $U(1)$-
transformationen $\psi\to e^{i\alpha}\psi$ svarer, skrevet ud i
real/imaginærdel, til $\varphi_1\to\varphi_1\cos\alpha-\varphi_2\sin\alpha$,
$\varphi_2\to\varphi_1\sin\alpha+\varphi_2\cos\alpha$ — en almindelig
2D-rotationsmatrix anvendt på $(\varphi_1,\varphi_2)$. Ladningen
$Q=\int\rho\,dx=2\mu\int(\varphi_1\dot\varphi_2-\varphi_2\dot\varphi_1)\,dx$
er dermed den kontinuerte (felt-)analog til banemomentet
$L_z=xp_y-yp_x$ — begge er Noether-ladningen for en rotation i et
($2$-dimensionalt) plan, blot for henholdsvis et diskret
partikelpar $(x,y)$ og et kontinuert feltpar $(\varphi_1,\varphi_2)$.

\section*{Kapitel 10: Hamiltonsk feltteori}
\addcontentsline{toc}{section}{Kapitel 10}

\paragraph{Opgave 1.} Hamiltons ligninger giver
$\dot\psi=\partial\mathcal{H}/\partial\pi=\pi^\ast/\mu$ og
$\dot\pi=-\partial\mathcal{H}/\partial\psi+d/dx(\partial\mathcal{H}/\partial\psi')
=0+\tau\psi^{\ast\prime\prime}$. Da $\pi=\mu\dot\psi^\ast$, giver
kombination af de to $\mu\ddot\psi^\ast=\tau\psi^{\ast\prime\prime}$,
dvs.\ $\ddot\psi^\ast=v^2\psi^{\ast\prime\prime}$ med $v^2=\tau/\mu$.
Tilsvarende for parret $(\psi^\ast,\pi^\ast)$ giver
$\ddot\psi=v^2\psi''$ — begge felter opfylder samme bølgeligning,
som forventet.

\paragraph{Opgave 2.} Analogt til \eqref{eq:Q-genererer}, men nu med
$\{\psi^\ast(x),\pi^\ast(y)\}=\delta(x-y)$ som det eneste bidragende
led i $Q=i\int(\psi\pi-\psi^\ast\pi^\ast)dx$ (med et minustegn foran
$\psi^\ast\pi^\ast$-leddet):
\begin{equation*}
  \{\psi^\ast(x),Q\} = -i\int\psi^\ast(y)\{\psi^\ast(x),\pi^\ast(y)\}\,dy
  = -i\,\psi^\ast(x).
\end{equation*}
Sammenholdt med $\delta\psi^\ast=-i\varepsilon\psi^\ast$ (fra
$U(1)$-transformationens infinitesimale form) bekræfter dette
$\delta\psi^\ast=\varepsilon\{\psi^\ast,Q\}$ — samme generator-mønster
som for $\psi$, nu for $\psi^\ast$.
